\documentclass{article}
\usepackage[a4paper,margin=2.4cm]{geometry}
\usepackage{graphicx}
\usepackage{float}
\usepackage{fancyvrb}
\fvset{frame=single,fontsize=\small,framesep=3mm}

\title{Labwork 6: Map}
\author{Pham Tien Nhat}
\date{\today}

\begin{document}

\maketitle

\section{Implementation}

The input image of 6a and 6b is loaded with \texttt{plt.imread()}.

\begin{Verbatim}
import math
import time
import numpy as np
import matplotlib.pyplot as plt

image = plt.imread('1.jpeg')
plt.imshow(image)
h, w, d = image.shape
\end{Verbatim}
\begin{figure}[H]
\centering
\includegraphics[width=0.45\textwidth]{../../image/lab6_input.png}
\caption{Input image of 6a and 6b.}
\label{figin}
\end{figure}

The input image is converted to gray with the 2D kernel before moving to binarization.

\begin{Verbatim}
@cuda.jit
def grayscale(src, dst):
    tidx = cuda.threadIdx.x + cuda.blockIdx.x * cuda.blockDim.x
    tidy = cuda.threadIdx.y + cuda.blockIdx.y * cuda.blockDim.y
    if tidx >= src.shape[1] or tidy >= src.shape[0]:
        return
    dst[tidy, tidx] = np.uint8((src[tidy, tidx, 0] + src[tidy, tidx, 1]
                                + src[tidy, tidx, 2]) / 3)

blockSize = (32, 32)
gridSize = (math.ceil(w / blockSize[0]), math.ceil(h / blockSize[1]))
devSrc = cuda.to_device(image)
devGray = cuda.device_array((h, w), np.uint8)
grayscale[gridSize, blockSize](devSrc, devGray)
plt.imshow(devGray.copy_to_host(), cmap='gray')
\end{Verbatim}
\begin{figure}[H]
\centering
\includegraphics[width=0.45\textwidth]{../../image/lab6_gray.png}
\caption{Gray image, input of 6a and 6b.}
\label{figgray}
\end{figure}

\textbf{6a: Binarization.} A pixel becomes 0 if its gray value is below the threshold, and 1 otherwise. I have tested with 3 different threshold values, where a lower threshold of 64 shows more white pixels, higher of 196 shows darker pixels.

\begin{Verbatim}
@cuda.jit
def binarize(src, dst, threshold):
    tidx = cuda.threadIdx.x + cuda.blockIdx.x * cuda.blockDim.x
    tidy = cuda.threadIdx.y + cuda.blockIdx.y * cuda.blockDim.y
    if tidx >= src.shape[1] or tidy >= src.shape[0]:
        return
    if src[tidy, tidx] < threshold:
        dst[tidy, tidx] = 0
    else:
        dst[tidy, tidx] = 1

devBin = cuda.device_array((h, w), np.uint8)
results = []
for threshold in [64, 128, 196]:
    binarize[gridSize, blockSize](devGray, devBin, threshold)
    plt.imshow(devBin.copy_to_host(), cmap='gray')
    plt.show()
    results.append(devBin.copy_to_host())

plt.figure(figsize=(15, 4))
for i, threshold in enumerate([64, 128, 196]):
    plt.subplot(1, 3, i + 1)
    plt.imshow(results[i], cmap='gray')
    plt.title(f"threshold = {threshold}")
plt.show()
\end{Verbatim}
\begin{figure}[H]
\centering
\includegraphics[width=\textwidth]{../../image/lab6_bin.png}
\caption{Combined output of 6a with $\tau = 64$, $128$ and $196$.}
\label{fig6a}
\end{figure}

\textbf{6b: Brightness control.} A value is added to every pixel and the result is limited to $[0, 255]$. 3 values are tested where a negative value of -80 shows dark image, and value of 120 shows much brighter image. 

\begin{Verbatim}
@cuda.jit
def brightness(src, dst, value):
    tidx = cuda.threadIdx.x + cuda.blockIdx.x * cuda.blockDim.x
    tidy = cuda.threadIdx.y + cuda.blockIdx.y * cuda.blockDim.y
    if tidx >= src.shape[1] or tidy >= src.shape[0]:
        return
    v = src[tidy, tidx] + value
    dst[tidy, tidx] = min(max(v, 0), 255)

devBright = cuda.device_array((h, w), np.uint8)
results = []
for value in [-80, 40, 120]:
    brightness[gridSize, blockSize](devGray, devBright, value)
    plt.imshow(devBright.copy_to_host(), cmap='gray',
               vmin=0, vmax=255)
    plt.show()
    results.append(devBright.copy_to_host())

plt.figure(figsize=(15, 4))
for i, value in enumerate([-80, 40, 120]):
    plt.subplot(1, 3, i + 1)
    plt.imshow(results[i], cmap='gray', vmin=0,
               vmax=255)
    plt.title(f"value = {value}")
plt.show()
\end{Verbatim}
\begin{figure}[H]
\centering
\includegraphics[width=\textwidth]{../../image/lab6_bright.png}
\caption{Combined output of 6b with value $= -80$, $40$ and $120$.}
\label{fig6b}
\end{figure}

\textbf{6c: Blending.} The two input images are cropped to a common size ($407 \times 415$).

\begin{Verbatim}
image2 = plt.imread('images.jpeg')
hb, wb = min(h, image2.shape[0]), min(w, image2.shape[1])
image1 = image[:hb, :wb].copy()
image2 = image2[:hb, :wb].copy()
plt.subplot(1, 2, 1)
plt.imshow(image1)
plt.subplot(1, 2, 2)
plt.imshow(image2)
plt.show()
\end{Verbatim}
\begin{figure}[H]
\centering
\includegraphics[width=0.8\textwidth]{../../image/lab6_in.png}
\caption{Input images for blending, cropped: \texttt{1.jpeg} (left) and \texttt{images.jpeg} (right).}
\label{figin2}
\end{figure}

2 images are combined using the formula: $Q = c \times \Phi_1 + (1 - c) \times \Phi_2$ on each RGB channel. A larger $c$ shows more of the first image, for example, c= 0.8 indicates that the output image will be 80% of the first input image and 20% of the second input image.

\begin{Verbatim}
@cuda.jit
def blend(src1, src2, dst, c):
    tidx = cuda.threadIdx.x + cuda.blockIdx.x * cuda.blockDim.x
    tidy = cuda.threadIdx.y + cuda.blockIdx.y * cuda.blockDim.y
    if tidx >= src1.shape[1] or tidy >= src1.shape[0]:
        return
    for k in range(3):
        dst[tidy, tidx, k] = (c * src1[tidy, tidx, k]
                              + (1 - c) * src2[tidy, tidx, k])

gridBlend = (math.ceil(wb / blockSize[0]), math.ceil(hb / blockSize[1]))
devSrc1 = cuda.to_device(image1)
devSrc2 = cuda.to_device(image2)
devBlend = cuda.device_array((hb, wb, 3), np.uint8)
results = []
for c in [0.2, 0.5, 0.8]:
    blend[gridBlend, blockSize](devSrc1, devSrc2, devBlend, c)
    plt.imshow(devBlend.copy_to_host())
    plt.show()
    results.append(devBlend.copy_to_host())

plt.figure(figsize=(15, 4))
for i, c in enumerate([0.2, 0.5, 0.8]):
    plt.subplot(1, 3, i + 1)
    plt.imshow(results[i])
    plt.title(f"c = {c}")
plt.show()
\end{Verbatim}
\begin{figure}[H]
\centering
\includegraphics[width=\textwidth]{../../image/lab6_blend.png}
\caption{Combined output of 6c with $c = 0.2$, $0.5$ and $0.8$.}
\label{fig6c}
\end{figure}

\section{Experimental Results}
4 2D block sizes are tested: $(4, 4)$, $(8, 8)$, $(16, 16)$ and $(32, 32)$, and the GPU time of each kernel is measured.

\begin{Verbatim}
block_sizes = [(4, 4), (8, 8), (16, 16), (32, 32)]
times_bin = []
times_bright = []
times_blend = []

for blockSize in block_sizes:
    gridSize = (math.ceil(w / blockSize[0]), math.ceil(h / blockSize[1]))

    start = time.time()
    binarize[gridSize, blockSize](devGray, devBin, threshold)
    cuda.synchronize()
    times_bin.append(time.time() - start)

    start = time.time()
    brightness[gridSize, blockSize](devGray, devBright, value)
    cuda.synchronize()
    times_bright.append(time.time() - start)

    start = time.time()
    gridBlend = (math.ceil(wb / blockSize[0]),
                 math.ceil(hb / blockSize[1]))
    blend[gridBlend, blockSize](devSrc1, devSrc2, devBlend, c)
    cuda.synchronize()
    times_blend.append(time.time() - start)

    print(f"Block size: {blockSize}")
    print(f"Grid size: {gridSize}")
    print(f"GPU time: {times_bin[-1]}, {times_bright[-1]},"
          f" {times_blend[-1]} seconds")

plt.figure(figsize=(8, 5))
plt.plot([str(b) for b in block_sizes], times_bin, marker='o',
         label='binarization')
plt.plot([str(b) for b in block_sizes], times_bright, marker='o',
         label='brightness')
plt.plot([str(b) for b in block_sizes], times_blend, marker='o',
         label='blending')
plt.xlabel("2D Block Size")
plt.ylabel("Time (seconds)")
plt.title("CUDA 2D Block Size vs GPU Execution Time")
plt.legend()
plt.grid(True)
plt.show()
\end{Verbatim}

Figure~\ref{fig6} shows the relationship between the 2D block size and the GPU time.

\begin{figure}[H]
\centering
\includegraphics[width=0.7\textwidth]{../../image/lab6.png}
\caption{2D block size vs GPU time.}
\label{fig6}
\end{figure}

The block of $(4, 4)$ is the slowest for all three kernels. Blending is the slowest kernel as it reads two images and writes three channels. 

\end{document}
